\chapter{Cuts, Mixing and Expansion}\label{ch:cheeger}
\begin{center}\itshape
The second eigenvalue of a regular graph says how cheaply the graph can be cut,\\
and how fast a random walk forgets where it started.
\end{center}

A connected graph can still be easy to split in two. The cheapest cut, relative to the piece it removes, is the Cheeger constant. The Laplacian's second eigenvalue bounds that cut, and the gap to the largest adjacency eigenvalue bounds how evenly the edges are spread. The graph spectra of the preceding chapter supply the eigenvalues. The expanders of the next chapter are the graphs for which the gap is as large as a tree allows.

\begin{lemma}[Rayleigh]\label{expansion:lem:rayleigh}
Let $M$ be a real symmetric matrix with eigenvalues $\nu_1\le\dots\le\nu_n$ and orthonormal eigenvectors $\psi_1,\dots,\psi_n$. If $x\ne0$ is orthogonal to
$\psi_1,\dots,\psi_{k-1}$, then $x^{\mathsf T}Mx\ge\nu_kx^{\mathsf T}x$, with equality for $x=\psi_k$. For the Laplacian $L$ of a graph, with eigenvalues $0=\mu_1\le\mu_2\le\dots\le\mu_n$,
\[
\mu_2=\min\Bigl\{\frac{x^{\mathsf T}Lx}{x^{\mathsf T}x}:x\ne0,\ x\perp\mathbf1\Bigr\}.
\]
\end{lemma}
\begin{proof}
Write $x=\sum_ic_i\psi_i$; then $x^{\mathsf T}x=\sum_ic_i^2$ and $x^{\mathsf T}Mx=\sum_i\nu_ic_i^2$. If $c_1=\dots=c_{k-1}=0$, this is at least $\nu_k\sum_ic_i^2$. For $L$, the first eigenvector is proportional to $\mathbf1$.
\end{proof}

\section{Cutting a graph: Cheeger's inequality}\label{spectral:sec:cheeger}

How well connected is a connected graph? A natural measure is the cheapest way to cut off a piece of it, relative to the size of the piece. Let $G$ be $d$-regular with $n$
vertices. For $S\subseteq V$, write $\bar S=V\setminus S$ and $e(S,\bar S)$ for the number of edges between $S$ and $\bar S$. The \emph{edge expansion}, or \emph{Cheeger constant}, of $G$ is
\[
h(G)=\min\Bigl\{\frac{e(S,\bar S)}{|S|}:S\ne\varnothing,\ |S|\le\frac n2\Bigr\}.
\]
Computing $h(G)$ exactly is NP-hard in general, a computational hardness property. The second eigenvalue gives it up to a square root, and its eigenvector gives a good cut.

\begin{theorem}[Cheeger's inequality: Dodziuk, 1984; Alon and Milman, 1985]\label{spectral:thm:cheeger}
For every $d$-regular graph,
\[
\frac{\mu_2}2\le h(G)\le\sqrt{2d\mu_2}.
\]
\end{theorem}
\begin{proof}
\emph{Lower bound.} Let $S$ achieve $h(G)$, $s=|S|\le\frac n2$, and $x=\mathbf1_S-\frac sn\mathbf1$, which is orthogonal to $\mathbf1$. An edge contributes $(x_u-x_v)^2=1$ if it crosses between
$S$ and $\bar S$ and $0$ otherwise, so $x^{\mathsf T}Lx=e(S,\bar S)$, while $x^{\mathsf T}x=s(1-\frac sn)^2+(n-s)\frac{s^2}{n^2}=\frac{s(n-s)}n$. By Lemma~\ref{expansion:lem:rayleigh},
\[
\mu_2\le\frac{n\,e(S,\bar S)}{s(n-s)}\le\frac{2e(S,\bar S)}s=2h(G),
\]
as $n-s\ge\frac n2$.

\emph{Upper bound.} Let $x$ be an eigenvector of $\mu_2$. It is orthogonal to $\mathbf1$ and non-zero, so it has positive and negative entries; replacing $x$ by $-x$ if necessary, the set
$V^+=\{v:x_v>0\}$ has at most $\frac n2$ elements. Let $y_v=\max(x_v,0)$.

First, $y^{\mathsf T}Ly\le\mu_2\,y^{\mathsf T}y$. For $v\in V^+$, $(Ly)_v=dy_v-\sum_{u\sim v}y_u\le dx_v-\sum_{u\sim v}x_u=(Lx)_v=\mu_2x_v=\mu_2y_v$, since $y_u\ge x_u$ and $y_v=x_v$. So
$y^{\mathsf T}Ly=\sum_{v\in V^+}y_v(Ly)_v\le\mu_2\sum_vy_v^2$.

Second, for $t\ge0$ let $S_t=\{v:y_v^2>t\}\subseteq V^+$, so $|S_t|\le\frac n2$ and $e(S_t,\bar S_t)\ge h(G)|S_t|$. An edge $uv$ crosses the cut $(S_t,\bar S_t)$ exactly for $t$ between $y_u^2$ and $y_v^2$,
an interval of length $|y_u^2-y_v^2|$. Integrating over $t$,
\[
\sum_{uv\in E}|y_u^2-y_v^2|=\int_0^\infty e(S_t,\bar S_t)\dd t\ge h(G)\int_0^\infty|S_t|\dd t=h(G)\sum_vy_v^2,
\]
since a vertex $v$ lies in $S_t$ for $t$ in an interval of length $y_v^2$.

Third, by the Cauchy--Schwarz inequality and $(a+b)^2\le2a^2+2b^2$,
\[
\begin{aligned}\sum_{uv\in E}|y_u^2-y_v^2|&=\sum_{uv\in E}|y_u-y_v|(y_u+y_v)\eqbr\le\sqrt{y^{\mathsf T}Ly}\,\sqrt{\textstyle\sum_{uv\in E}2(y_u^2+y_v^2)}=\sqrt{y^{\mathsf T}Ly}\,\sqrt{2d\textstyle\sum_vy_v^2},\end{aligned}
\]
as each vertex is in $d$ edges. Combining the three steps, $h(G)\sum_vy_v^2\le\sqrt{2d\mu_2}\sum_vy_v^2$, and $y\neq0$.
\end{proof}

The proof is an algorithm, called \emph{spectral partitioning}. The integral inequality $\int_0^\infty e(S_t,\bar S_t)\dd t\le\sqrt{2d\mu_2}\int_0^\infty|S_t|\dd t$ shows that some $t$ satisfies
$e(S_t,\bar S_t)\le\sqrt{2d\mu_2}\,|S_t|$, and the sets $S_t$ are the $n$ or fewer \emph{sweep sets}: sort the vertices by the value of the Fiedler vector and try every threshold. In 2000 Jianbo Shi and
Jitendra Malik used a normalised version of this idea to segment images, the pixels being vertices joined according to their similarity, and \emph{spectral clustering} became a standard tool of
data analysis.

Both inequalities are tight up to constants. For the hypercube $Q_n$, $\mu_2=2$ and $h(Q_n)=1$ (cut along a coordinate; that nothing is better is the edge-isoperimetric inequality, a classical
exercise of combinatorics), so the left side is attained. For the cycle, $\mu_2\approx\frac{4\pi^2}{n^2}$ while $h(C_n)\approx\frac4n$, so the square root on the right is necessary (Exercise~\ref{spectral:ex:cheegercycle}).

\begin{figure}[ht]
\centering
\fig{spectral-ch09-fiedler}
\caption{Top: a graph made of two dense clusters of $11$ vertices joined by three edges (rust). Bottom: the values of its Fiedler vector $\varphi_2$, one dot per vertex. The best
sweep set (dashed threshold) is exactly one cluster (charcoal): the cut crosses only the three bridges.}
\label{spectral:fig:fiedler}
\end{figure}

\section{Expanders}

For a $d$-regular graph, let $d=\lambda_1\ge\lambda_2\ge\dots\ge\lambda_n$ be the eigenvalues of $A$ (so $\mu_i=d-\lambda_i$), and $\lambda=\max(|\lambda_2|,|\lambda_n|)$, the largest absolute value of an
eigenvalue other than $d$. A $d$-regular graph on $n$ vertices with this $\lambda$ is called an \emph{$(n,d,\lambda)$-graph}. Cheeger's inequality says that $h(G)\ge\frac{d-\lambda_2}2$: if $\lambda$ is small compared
with $d$, every set of at most half the vertices has many edges leaving it. A family of $d$-regular graphs with $n\to\infty$ and $\lambda\le(1-\varepsilon)d$ is a family of \emph{expanders}. The
next lemma says that in such a graph, edges are spread as if at random.

\begin{lemma}[expander mixing lemma, Alon and Chung, 1988]\label{spectral:lem:mixing}
Let $G$ be an $(n,d,\lambda)$-graph and $S,T\subseteq V$. Let $e(S,T)$ be the number of pairs $(s,t)\in S\times T$ with $s\sim t$. Then
\[
\Bigl|e(S,T)-\frac{d|S||T|}n\Bigr|\le\lambda\sqrt{|S||T|\Bigl(1-\frac{|S|}n\Bigr)\Bigl(1-\frac{|T|}n\Bigr)}\le\lambda\sqrt{|S||T|}.
\]
\end{lemma}
\begin{proof}
Write $\mathbf1_S=\frac{|S|}n\mathbf1+x$ and $\mathbf1_T=\frac{|T|}n\mathbf1+y$, with $x,y\perp\mathbf1$. Then $e(S,T)=\mathbf1_S^{\mathsf T}A\mathbf1_T$, and since $A\mathbf1=d\mathbf1$ and $A$ maps $\mathbf1^\perp$ to itself,
\[
e(S,T)=\frac{|S||T|}{n^2}\,\mathbf1^{\mathsf T}A\mathbf1+x^{\mathsf T}Ay=\frac{d|S||T|}n+x^{\mathsf T}Ay.
\]
Expand $x=\sum_{i\ge2}c_i\psi_i$ and $y=\sum_{i\ge2}c'_i\psi_i$ in an orthonormal eigenbasis of $A$ with $\psi_1\propto\mathbf1$: $|x^{\mathsf T}Ay|=|\sum_{i\ge2}\lambda_ic_ic'_i|\le\lambda\|x\|\|y\|$ by the Cauchy--Schwarz
inequality. Finally $\|x\|^2=\|\mathbf1_S\|^2-\frac{|S|^2}n=|S|(1-\frac{|S|}n)$, and likewise for $y$.
\end{proof}

For instance an independent set $S$ has $e(S,S)=0$, so $\frac{d|S|^2}n\le\lambda|S|(1-\frac{|S|}n)$, that is $|S|\le\frac{n\lambda}{d+\lambda}$ (compare Hoffman's bound, Volume~8, Theorem~9.6). The lemma also
bounds the diameter (Exercise~\ref{spectral:ex:diameter}) and shows that random walks forget their starting point exponentially fast (Exercise~\ref{spectral:ex:walkmixing}). Good expanders want $\lambda$ as small as
possible. How small can it be?

\begin{theorem}[Alon and Boppana, 1986; Nilli's proof, 1991]\label{spectral:thm:alonboppana}
Let $G$ be $d$-regular, $d\ge2$, with two edges $\{v_1,v_2\}$ and $\{u_1,u_2\}$ whose endpoints are at distance at least $2k+2$ from each other. Then
\[
\lambda_2\ge2\sqrt{d-1}-\frac{2\sqrt{d-1}-1}{k+1}.
\]
\end{theorem}
\begin{proof}
We build a test vector for Lemma~\ref{expansion:lem:rayleigh} that looks, near each of the two edges, like the eigenvector of an infinite $d$-regular tree. Let $V_i$ be the set of vertices at distance
exactly $i$ from $\{v_1,v_2\}$, for $0\le i\le k$, and $a_i=(d-1)^{-i/2}$. Define $f_x=a_i$ for $x\in V_i$, $i\le k$, and $f_x=0$ elsewhere.

Every vertex of $V_i$ with $i\ge1$ has a neighbour in $V_{i-1}$, and each of $v_1,v_2$ has the other as a neighbour in $V_0$; so each vertex of $V_i$ has at most $d-1$ neighbours in $V_{i+1}$, and
$|V_{i+1}|\le(d-1)|V_i|$. Edges inside some $V_i$, or outside $V_0\cup\dots\cup V_{k+1}$, contribute nothing to $f^{\mathsf T}Lf=\sum_{xy\in E}(f_x-f_y)^2$. Edges between $V_i$ and $V_{i+1}$ number at most
$(d-1)|V_i|$, so
\[
f^{\mathsf T}Lf\le\sum_{i=0}^{k-1}(d-1)|V_i|(a_i-a_{i+1})^2+(d-1)|V_k|a_k^2.
\]
Now $(d-1)(a_i-a_{i+1})^2=a_i^2(\sqrt{d-1}-1)^2=a_i^2(d-2\sqrt{d-1})$, and $(d-1)a_k^2=a_k^2(d-2\sqrt{d-1})+a_k^2(2\sqrt{d-1}-1)$. Hence
\[
f^{\mathsf T}Lf\le(d-2\sqrt{d-1})\sum_{i=0}^k|V_i|a_i^2+(2\sqrt{d-1}-1)|V_k|a_k^2.
\]
The numbers $|V_i|a_i^2$ do not increase with $i$, since $|V_{i+1}|a_{i+1}^2\le(d-1)|V_i|\frac{a_i^2}{d-1}$; so the last one is at most their average, $|V_k|a_k^2\le\frac1{k+1}\sum_i|V_i|a_i^2=\frac{f^{\mathsf T}f}{k+1}$. With
$\theta=d-2\sqrt{d-1}+\frac{2\sqrt{d-1}-1}{k+1}$ we get $f^{\mathsf T}Lf\le\theta f^{\mathsf T}f$.

Define $g$ in the same way around $\{u_1,u_2\}$; also $g^{\mathsf T}Lg\le\theta g^{\mathsf T}g$. The supports of $f$ and $g$ are at distance at least $2k+2-2k=2$, so they are disjoint and no edge joins them.
Choose $c$ with $h=f-cg\perp\mathbf1$ (possible as $\sum g>0$). Then $h\ne0$, and since no edge meets both supports, $h^{\mathsf T}Lh=f^{\mathsf T}Lf+c^2g^{\mathsf T}Lg\le\theta(f^{\mathsf T}f+c^2g^{\mathsf T}g)=\theta h^{\mathsf T}h$. By
Lemma~\ref{expansion:lem:rayleigh}, $\mu_2\le\theta$, that is $\lambda_2=d-\mu_2\ge d-\theta$.
\end{proof}

A $d$-regular graph with $n$ vertices has at most $2(1+(d-1)+\dots+(d-1)^r)$ vertices within distance $r$ of a given edge, so when $n\to\infty$ with $d$ fixed, there are edges at distance $2k+2$ for larger
and larger $k$: $\lambda_2\ge2\sqrt{d-1}-o(1)$. The number $2\sqrt{d-1}$ is the spectral radius of the infinite $d$-regular tree, which every large $d$-regular graph resembles locally. Graphs with
$\lambda\le2\sqrt{d-1}$ are called \emph{Ramanujan graphs}: Lubotzky, Phillips and Sarnak (1988), and independently Margulis, built them for $d-1$ prime using Ramanujan's conjecture on the
coefficients of modular forms (proved by Deligne), and Marcus, Spielman and Srivastava (2013) proved that bipartite Ramanujan graphs exist for every degree, with a new method of ``interlacing
families'' of polynomials. The next chapter tells this story. And Joel Friedman proved in 2008 that a \emph{random} $d$-regular graph has $\lambda\le2\sqrt{d-1}+\varepsilon$ with probability tending to $1$: almost all
regular graphs are almost optimal expanders, while building one explicitly is hard.

\begin{beauty}
Two numbers attached to a matrix, the second smallest Laplacian eigenvalue and the gap to the spectral radius of a tree, answer two questions that seem to have nothing to do with
linear algebra: how cheaply a network can be cut, and how fast a walk forgets its start. The count of the spanning trees, the product of the non-zero eigenvalues, is proved in Volume~12.
\end{beauty}

\section{Exercises}
\begin{exercise}[$\star$]\label{spectral:ex:cheegercycle}
Show that $h(C_n)=\frac2{\lfloor n/2\rfloor}$ and $\mu_2(C_n)=2-2\cos\frac{2\pi}n$. Conclude that in Theorem~\ref{spectral:thm:cheeger} the upper bound cannot be replaced by $c\,\mu_2$ for any constant $c$.
\end{exercise}
\begin{exercise}[$\star\star$]\label{spectral:ex:diameter}
Let $G$ be a connected $(n,d,\lambda)$-graph with $\lambda<d$. Show that $(A^t)_{uv}\ge\frac{d^t}n-\lambda^t$ for all vertices $u,v$, and deduce that the diameter of $G$ is at most
$\bigl\lfloor\frac{\ln n}{\ln(d/\lambda)}\bigr\rfloor+1$.
\end{exercise}
\begin{exercise}[$\star$]\label{spectral:ex:walkmixing}
A random walk on a $d$-regular graph moves from its current vertex to a uniform random neighbour. Let $p_t$ be the distribution after $t$ steps and $u=\frac1n\mathbf1$. Show that $p_{t+1}=\frac1dAp_t$
and $\|p_t-u\|\le(\frac\lambda d)^t$ for an $(n,d,\lambda)$-graph.
\end{exercise}
\begin{exercise}[$\star$]\label{spectral:ex:petersenalpha}
Use Lemma~\ref{spectral:lem:mixing} to show that the Petersen graph has no independent set of $5$ vertices. Find one of $4$.
\end{exercise}
\begin{exercise}[$\star\star$]\label{spectral:ex:fiedlerpath}
Show that the eigenvalues of the Laplacian of the path $P_n$ are $2-2\cos\frac{j\pi}n$, $j=0,\dots,n-1$, with eigenvectors $x_v=\cos\frac{j\pi(2v-1)}{2n}$, $v=1,\dots,n$. What cut does the sweep of
the Fiedler vector find?
\end{exercise}
\begin{exercise}[$\star$]\label{spectral:ex:moorebound}
(a) Show that at most $2\bigl(1+(d-1)+\dots+(d-1)^r\bigr)$ vertices of a $d$-regular graph are within distance $r$ of a given edge, with equality when the graph has girth at least $2r+2$.
(b) Deduce that a $d$-regular graph with more than $2\bigl(1+(d-1)+\dots+(d-1)^{2k+2}\bigr)$ vertices has two edges at distance at least $2k+2$, so that Theorem~\ref{spectral:thm:alonboppana} applies.
\end{exercise}
\section{Hints}
\begin{hint}{spectral:ex:cheegercycle}
A set of $s$ vertices of a cycle has at least $2$ boundary edges, exactly $2$ for an arc.
\end{hint}
\begin{hint}{spectral:ex:diameter}
Use the spectral decomposition $A^t=\sum_i\lambda_i^t\psi_i\psi_i^{\mathsf T}$ and $\sum_i\psi_i(u)^2=1$.
\end{hint}
\begin{hint}{spectral:ex:walkmixing}
$p_t-u$ is orthogonal to $\mathbf1$.
\end{hint}
\begin{hint}{spectral:ex:petersenalpha}
$\lambda=2$ for the Petersen graph.
\end{hint}
\begin{hint}{spectral:ex:fiedlerpath}
Check $(Lx)_v=\mu x_v$ at the inner vertices with $\cos(a-b)+\cos(a+b)=2\cos a\cos b$, then at the ends.
\end{hint}
\begin{hint}{spectral:ex:moorebound}
Count the vertices at distance exactly $i$, as in the proof of Theorem~\ref{spectral:thm:alonboppana}.
\end{hint}
\section{Solutions}
\begin{solution}{spectral:ex:cheegercycle}
\idea Cutting an arc of the cycle costs two edges; the eigenvalue is tiny because a slowly varying cosine has small energy.
\writeup A non-empty proper subset $S$ of the cycle is a union of $r\ge1$ arcs, with $2r\ge2$ boundary edges, so $\frac{e(S,\bar S)}{|S|}\ge\frac2{|S|}\ge\frac2{\lfloor n/2\rfloor}$, with equality for an arc of $\lfloor\frac n2\rfloor$ vertices. The
eigenvalues of the cycle are $2-2\cos\frac{2\pi j}n$, the smallest non-zero one for $j=1$. As $n\to\infty$, $\mu_2\sim\frac{4\pi^2}{n^2}$ while $h(C_n)\sim\frac4n$: the ratio $\frac h{\mu_2}\sim\frac n{\pi^2}$ is unbounded, while
$\frac h{\sqrt{2d\mu_2}}\to\frac1\pi$ ($d=2$). So no bound $h\le c\mu_2$ holds, and the square root is the right order.
\end{solution}
\begin{solution}{spectral:ex:diameter}
\idea Count walks: the main term $\frac{d^t}n$ from the eigenvalue $d$ beats the error from the others once $t$ is large, and then every pair of vertices is joined by a walk of length $t$.
\writeup With an orthonormal eigenbasis, $\psi_1=\frac1{\sqrt n}\mathbf1$, $(A^t)_{uv}=\sum_i\lambda_i^t\psi_i(u)\psi_i(v)=\frac{d^t}n+\sum_{i\ge2}\lambda_i^t\psi_i(u)\psi_i(v)$. The rows of the orthogonal matrix $(\psi_i(u))$ have norm $1$, so
by Cauchy--Schwarz $\sum_{i\ge2}|\psi_i(u)\psi_i(v)|\le1$, and $(A^t)_{uv}\ge\frac{d^t}n-\lambda^t$. If $(\frac d\lambda)^t>n$, this is positive, so there is a walk of length $t$ from $u$ to $v$, and the distance is
at most $t$. Note that $\lambda>0$, as the eigenvalues add up to $\operatorname{tr}A=0$. The integer $t=\lfloor\frac{\ln n}{\ln(d/\lambda)}\rfloor+1$ exceeds $\frac{\ln n}{\ln(d/\lambda)}$, so $(\frac d\lambda)^t>n$ and the diameter is at most $t$. Expanders have logarithmic diameter.
\end{solution}
\begin{solution}{spectral:ex:walkmixing}
\idea The uniform distribution is the eigenvector of the top eigenvalue, and everything orthogonal to it shrinks by a factor $\frac\lambda d$ per step.
\writeup $p_{t+1}(v)=\sum_{u\sim v}\frac1dp_t(u)$, that is $p_{t+1}=\frac1dAp_t$ ($A$ is symmetric). Also $\frac1dAu=u$. So $p_{t+1}-u=\frac1dA(p_t-u)$, and $p_t-u$ has total mass $0$, so it lies in $\mathbf1^\perp$, where
$\frac1dA$ has eigenvalues of absolute value at most $\frac\lambda d$: $\|p_{t+1}-u\|\le\frac\lambda d\|p_t-u\|$. By induction $\|p_t-u\|\le(\frac\lambda d)^t\|p_0-u\|$, and $\|p_0-u\|^2=1-\frac1n\le1$ for a walk started at a vertex. A bipartite
graph has $\lambda_n=-d$, and the walk indeed never mixes: it alternates between the sides.
\end{solution}
\begin{solution}{spectral:ex:petersenalpha}
\idea The mixing lemma with $S=T$ independent.
\writeup The Petersen graph is a $(10,3,2)$-graph (eigenvalues $3,1,-2$). An independent set has $|S|\le\frac{n\lambda}{d+\lambda}=\frac{20}5=4$. In the usual drawing as the $2$-subsets of $\{1,\dots,5\}$, adjacent when disjoint,
the four $2$-sets containing $1$ pairwise intersect, so they are independent. Here the mixing bound is as good as Hoffman's.
\end{solution}
\begin{solution}{spectral:ex:fiedlerpath}
\idea Cosines satisfy the three-term recurrence, and the half-integer shift makes the ends behave like reflecting walls.
\writeup Let $\theta=\frac{j\pi}n$ and $x_v=\cos((v-\frac12)\theta)$. For $1<v<n$, $(Lx)_v=2x_v-x_{v-1}-x_{v+1}=2x_v-2\cos\theta\,x_v=(2-2\cos\theta)x_v$. At $v=1$, $(Lx)_1=x_1-x_2$, and the formula for the
inner vertices would give $2x_1-x_0-x_2$; they agree because $x_0=\cos(-\frac\theta2)=x_1$. Similarly, at $v=n$ we need $x_{n+1}=x_n$, and indeed $x_{n+1}=\cos(j\pi+\frac\theta2)$ and $x_n=\cos(j\pi-\frac\theta2)$ are both $(-1)^j\cos\frac\theta2$. Each $x$
is non-zero ($x_1=\cos\frac\theta2>0$), and the $n$ eigenvalues, for $j=0,\dots,n-1$, are distinct, so these are all the eigenvalues. The Fiedler vector ($j=1$) is a decreasing cosine; its sweep sets are the initial segments of the path, and the best is the
half path, cut in the middle.
\end{solution}
\begin{solution}{spectral:ex:moorebound}
\idea Grow balls around the edge: each new layer is at most $d-1$ times the previous one.
\writeup (a) Let $V_i$ be the set of vertices at distance $i$ from the edge $\{v_1,v_2\}$. As in the proof of Theorem~\ref{spectral:thm:alonboppana}, $|V_0|=2$ and $|V_{i+1}|\le(d-1)|V_i|$, so $|V_i|\le2(d-1)^i$, and summing over
$i\le r$ gives the bound. If the girth is at least $2r+2$, the vertices within distance $r$ of $v_1$ on the side away from $v_2$, and those of $v_2$, form two disjoint trees in which every vertex at
distance $i<r$ has $d-1$ children (a coincidence or an extra edge would close a cycle of length at most $2r+1$), so all inequalities are equalities.
(b) By (a) with $r=2k+2$, some vertex $x$ is at distance at least $2k+3$ from $\{v_1,v_2\}$. Take any edge $\{x,x'\}$: $x'$ is at distance at least $2k+2$. So the two edges are at distance at least
$2k+2$, and $\lambda_2\ge2\sqrt{d-1}-\frac{2\sqrt{d-1}-1}{k+1}$. For fixed $d$ and $n\to\infty$, $k$ can be taken of order $\log_{d-1}n$, which proves $\lambda_2\ge2\sqrt{d-1}-O(\frac1{\log n})$.
\end{solution}
